\documentclass[12pt,
				 a4paper,
				%landscape,
				%headings=Big,
				numbers=endperiod
				 ]{scrartcl}
% !TeX root = f61_longreport_schmitt_kleinbek.tex
\usepackage{polyglossia}				%Sprache
\usepackage{amsmath}					%Matheumgebung
\usepackage{amsopn}					%Matheoperatoren
\usepackage{amssymb}					%Mathesymbole
\usepackage{amsthm}					%Mathetheorem
\usepackage{array}						%Tabellenumgebung
\usepackage{booktabs}					%Tabellen
\usepackage[style=english]{csquotes}	%Anführungszeichen
\usepackage[includeheadfoot,
			top=1cm,
			bottom=1.5cm,
			left=2cm,
			right=2cm,
			]{geometry}
\usepackage{graphicx}					%Graphik
\usepackage{mathdots}					%Punkte
\usepackage{mathtools}				%Bugfix ams
\usepackage{microtype}				%Makrotypographie
\usepackage{nicefrac}
\usepackage{physics}
\usepackage{scrhack}
\usepackage[headsepline, automark]
			{scrlayer-scrpage}			%Kopfzeile
\usepackage[separate-uncertainty=true,
			quotient-mode=fraction,
			per-mode=symbol,
			list-final-separator={,},
			]{siunitx}					%Einheiten
\usepackage{subcaption}				%Gleitumgebung Graphik
\usepackage{upgreek}					%Griechische Buchstaben
\usepackage{xcolor}					%Farben
\usepackage{xltxtra}					%fontec

\setmainlanguage{german}
\setmainfont{Linux Libertine O}
\setsansfont{Linux Biolinum O}
%\setmonofont{DejaVu Sans Mono}

\ihead{F70 Mechanik und Vakuum}
\ohead{\headmark}
\chead{}
\pagestyle{scrheadings}

\setcounter{tocdepth}{1}
\setcounter{secnumdepth}{1}
\setlength{\parindent}{0pt}

\usepackage[bookmarksopenlevel=section,
			linkcolor=blue,
			colorlinks=true,
			urlcolor=blue,
			citecolor=blue]
			{hyperref}		%Verweise, muss am Ende stehen

\titlehead{\begin{minipage}{.5\textwidth}
			\begin{flushleft}
			\includegraphics[scale=.065]{logo_color.jpg}
			\end{flushleft}
 	 		\end{minipage}
			\begin{minipage}{.5\textwidth}
			\begin{flushright}
			\includegraphics[scale=.8]{logo_fp.png}
			\end{flushright}
 	 		\end{minipage}}
\title{F70 Mechanik und Vakuum}
\subtitle{Kritische Würdigung}
\author{Nils Schmitt}
\publishers{Tutor: Bätzner, Rolf}
\date{Durchgeführt im September 2018}

\begin{document}
\maketitle
In diesem Versuch konnten wir die Arbeit eines Vakuumphysikers mittels zwei verschiedenen Vakuumumpen kennenlernen.\\

Dabei ließ sich an der Drehschieberpumpe der sogenannte Gasballast austesten um einen Basisdruck im Feinvakuumgebiet zu erreichen.
Mit der Wasserschale in der Vakuumkammer konnten wir so die Zustandsänderung von Wasser im Phasendiagramm nachvollziehen.
Auch wenn sich der Tripelpunkt leider nicht beobachten lies.\\
Der Basisdruck der Turbomolekularpumpe (TMP) ging bei uns bis unter \SI{e-6}{\milli\bar}.
Dies konnten wir aber nur über die lange Laufzeit erreichen, was jedoch schneller durch eine größere Öffnung zum Rezipienten gegangen wäre.
So konnten wir ein effektives Saugvermögen der TMP von $S_{\text{eff}}=\SI{65.6(11)}{\liter\per\second}$ bestimmen.
Der Fehler der Messung kommt hauptsächlich aus dem Zeitfehler, welcher durch längere Teststrecken deutlich hätte verkleinert werden können.
Außerdem waren gerade die größeren Kapillare auch nicht ganz sauber zu bekommen, sodass der Durchmesser an manchen Stellen kleiner war.
Beim Kolbenmesser sollte dann noch angemerkt werden, dass wir den vom herrschenden Luftdruck sicher verschiedenen Normdruck verwendet haben, wodurch auch hier eine Korrektur möglich ist.\\

Die anschließende Leitwertmessung von Rohr, Blende und der Kombination konnten wir die Kirchhoffschen Regeln ür die Reihenschaltung bestätigen.
Genauso weicht der Leitwert der Blende nicht signifikant von dem theoretischen Wert ab.
Der theoretische Leitwert des Rohs hingegen hängt bei viskoser Strömung vom Druck ab, während dies bei der Molekularströmung nicht der Fall ist.
In dem Bereich, in dem unsere Messung liegt, treten aber beide Strömungen auf, sodass unser Messwert zwischen den Beiden erwarteten liegt.\\

Insgesamt muss man über den Versuch sagen, dass weit mehr als nur das Wissen über ein bisschen Vakuum vermittelt wurde.
So wurde zu Beginn die Abfrage durch viel Fragen und Antworten von Allgemeinwissen ergänzt.
Außerdem konnten wir bei Umbauten auch selbst etwas handwerkliche Fähigkeiten zeigen.
Daran merkt man auch, dass der Versuch alles andere als eine \enquote{Black Box} ist.
Die verschiedenen Pumpensysteme wurden schließlich auch noch in verschiedenen Graphiken dargestellt.
Anzumerken ist außerdem, dass es keine Selbstverständlichkeit ist, dass der Tutor die ganze Zeit anwesend ist, was aber hier zum wesentlichen Teil dem Lernerfolg beigetragen hat.
Gerade der eingeschobene Teil mit dem Glasblasen hat mir sehr große Freude gemacht.\\
So bleibt mir also nicht weiteres als zu sagen, dass sich dieser Versuch in der ganzen Länge von zwei Tagen vollständig lohnt, um auch mal ein bisschen über den Tellerrand der Physik herauszuschauen.
\end{document}